%!TeX program = lualatex

\documentclass[oneside]{academica}
\usetikzlibrary{automata,positioning,trees,patterns} % for automata and trees
\usepackage{diagbox} % for boxes
\usepackage{scalerel} % for scaling symbols
\usepackage{stmaryrd} % for llbracket

\setlength{\parindent}{0pt}

\hyphenation{ta-bleau}
\hyphenation{ta-bleaux}

\newcommand{\wlogt}{w.l.o.g.\xspace}
\newcommand{\Wlogt}{W.l.o.g.\xspace}

\newcommand{\DBA}{\tuple{Q, A, \d, q_0, F}}
\newcommand{\NBA}{\tuple{Q, A, \D, q_0, F}}
\newcommand{\DMA}{\tuple{Q, A, \d, q_0, \F}}
\newcommand{\NMA}{\tuple{Q, A, \D, q_0, \F}}
\newcommand{\DRA}{\tuple{Q, A, \d, q_0, \O}}
\newcommand{\NRA}{\tuple{Q, A, \D, q_0, \O}}

\newcommand{\Aw}{A^\w}

%\newcommand{\eNFA}{{{$\varepsilon$}-NFA\xspace}}
\newcommand{\eNFA}{{e-NFA\xspace}}

\newcommand{\Buchi}{B\"{u}chi\xspace}
\newcommand{\wword}{$\w$-word\xspace}
\newcommand{\wwords}{$\w$-words\xspace}
\newcommand{\wclo}{$\w$-closure\xspace}
\newcommand{\wlan}{$\w$-language\xspace}
\newcommand{\wlans}{$\w$-languages\xspace}
\newcommand{\wreg}{$\w$-regular\xspace}

\newcommand{\SA}{$\text{S1S}_A$\xspace}

\renewcommand{\succ}{\text{succ}}
\newcommand{\succxy}{\succ(X,Y)}
\newcommand{\tsuc}{$\t$-successor\xspace}

\newcommand{\sing}[1]{\text{sing}(#1)}

\renewcommand{\inf}[1]{\text{Inf}(#1)}
\newcommand{\infs}{\inf{\sigma}}

\newcommand{\eqn}{\equiv_{n}}
\newcommand{\eqnn}{\equiv_{(n,1)}}

\newcommand{\px}{\p(x)}

\newcommand{\Twa}{T^\w_A}

\newcommand{\AP}{{\mathcal{AP}}}
\newcommand{\sAP}{(2^{\AP})^\w}

\newcommand{\wY}{{\tilde{Y}}}
\newcommand{\state}[1]{\text{state}(#1)}

\newcommand{\bgraph}{\B_{(\P,\p)}}

\newcommand{\mcirc}{{\mathbin{\scalerel*{\bigcirc}{gX}}}}
\newcommand{\msquare}{{\mathbin{\scalerel*{\square}{gX}}}}

\renewcommand{\emptyset}{\varnothing}

\newcommand{\XX}{\mathrm{X}}
\newcommand{\UU}{\mathrm{U}}
\newcommand{\RRR}{\mathrm{R}}
\newcommand{\EE}{\mathrm{E}}
\renewcommand{\AA}{\mathrm{A}}
\newcommand{\FF}{\mathrm{F}}
\newcommand{\GG}{\mathrm{G}}
\newcommand{\EX}{\mathrm{EX}}
\newcommand{\AX}{\mathrm{AX}}
\newcommand{\EU}{\mathrm{EU}}
\newcommand{\AU}{\mathrm{AU}}
\newcommand{\EF}{\mathrm{EF}}
\newcommand{\AF}{\mathrm{AF}}
\newcommand{\EG}{\mathrm{EG}}
\newcommand{\AG}{\mathrm{AG}}
\newcommand{\FG}{\mathrm{FG}}
\newcommand{\GF}{\mathrm{GF}}

\newcommand{\ol}[1]{\,\overline{\!{#1}}}
\newcommand{\ul}[1]{\,\underline{{#1}\!}}
%\newcommand{\oa}[1]{\overrightarrow{{#1}\!}}

\newcommand{\lLk}{{}_l L_k}
\newcommand{\kloop}[1]{\llbracket {#1} \rrbracket_k}
\newcommand{\kloopl}[1]{{}_l \llbracket {#1} \rrbracket_k}
\newcommand{\kloopil}[1]{{}_l \llbracket {#1} \rrbracket^i_k}
\newcommand{\kloopsl}[1]{{}_l \llbracket {#1} \rrbracket^{\succ(i)}_k}

\newcommand{\var}{\text{var}}
\newcommand{\val}{\text{value}}
\newcommand{\low}{\text{low}}
\newcommand{\high}{\text{high}}

\newcommand{\attr}{\text{Attr}}


\overtitle{oral exam questions for}
\title{Verification and Validation}
\author{Transcribed from handwritten notes}

\begin{document}

\frontmatter

\maketitle

\tableofcontents

\mainmatter

%======================================================================
\chapter{How to read these notes}
%======================================================================

These notes reconstruct, in a clean and structured form, the questions
asked during a session of \emph{oral} exams of the Verification and
Validation course. They were taken live, by hand, in a mixture of
English and Italian, and are deliberately telegraphic; what follows is a
faithful transcription together with a topic-oriented study companion.

\section*{Source and method}

The raw material is a set of $18$ handwritten pages
(\texttt{domande.pdf}). Each oral exam there begins with the topic
\emph{chosen by the student} (written \quoted{choice} / \quoted{argomento
a scelta}), followed by the examiner's follow-up questions, and usually
ends with a grade. Wherever the handwriting was incomplete, ambiguous,
or simply wrong, the gap has been filled using the course notes
(\texttt{notes\_restructured.tex}) and standard knowledge of the field.
Every such decision is logged in the companion file
\texttt{ambiguities.md}.

\section*{Conventions}

\begin{itemize}
  \item \textbf{Two parts.} \autoref{cha:bysession} is the
    \emph{faithful} transcription --- ``what was asked'' in each single
    exam, in the original order. \autoref{cha:bytopic} is the
    \emph{derived}, study-oriented reorganisation of the same questions
    by topic, with short model answers and pointers back to the notes.
  \item \textbf{Square brackets} \texttt{[\,\ldots\,]} mark text that is
    \emph{not} in the original scribbles but reconstructed by me
    (intended answer, missing operator, the property a question refers
    to, etc.).
  \item A trailing \quoted{(?)} marks a low-confidence reading.
  \item \textbf{Grades} follow the Italian scale: $18$ is the pass mark,
    $30$ the maximum, and $30L$ means \emph{30 e lode} (30 cum laude).
    Verdicts such as \quoted{failed}, \quoted{left blank}
    (\textit{lasciato}) or \quoted{knew everything} (\textit{saputo
    tutto}) are reported as written.
  \item \textbf{LTL} operators are written in plain italics
    ($X,U,F,G$); \textbf{CTL} operators use the upright macros
    ($\EX,\AX,\EF,\AF,\EG,\AG,\AU,\EU$), exactly as in the course notes.
\end{itemize}

\begin{note}
  Several exams in the source carry no explicit \quoted{choice} header
  (the chosen topic was not written down) and a few carry no grade.
  These omissions are reported as \quoted{not recorded}; they are gaps
  in the original notes, not in the transcription.
\end{note}

%======================================================================
\chapter{What was asked --- exam by exam}\label{cha:bysession}
%======================================================================

This chapter follows the original page order. Fifteen distinct sittings
could be separated (each starts with a student-chosen topic). The table
below is an index; the grade column reproduces the mark written at the
end of each exam.

\begin{center}
  \small
  \begin{tabular}{@{}clcc@{}}
    \textbf{\#} & \textbf{Chosen topic / first question} & \textbf{Lang.} & \textbf{Grade} \\
    \hline
    1  & S1S $\equiv$ NBA --- the easy direction of the proof          & en & $30L$ \\
    2  & Model checking: definition                                    & en & $19$ \\
    3  & Overview of all the automaton types                           & en & --- (Q3 wrong) \\
    4  & Learning automata                                             & en & $18$ \\
    5  & $\approx_{\A}$ saturates / finite index / partitions          & en & not rec. \\
    6  & LTL                                                           & en & Q5 failed \\
    7  & [not recorded] --- starts at \quoted{finitely many $a$}       & en & $23$ \\
    8  & [not recorded] --- DBA / S1S / OBDD                           & en & not rec. \\
    9  & DBA                                                          & en & not rec. \\
    10 & LTL+P succinctness                                            & it & $29$ \\
    11 & Closure of BA under intersection                              & it & not rec. \\
    12 & S1S $\equiv$ $\w$-regular (Büchi's theorem)                   & it & $30L$ \\
    13 & Muller automata $\equiv$ Boolean combinations                 & it & $28$ \\
    14 & Closure of $\w$-regular languages under union                 & it & left blank \\
    15 & Learning automata                                             & it & $30L$ \\
  \end{tabular}
\end{center}

%----------------------------------------------------------------------
\section{Exam 1 --- S1S $\equiv$ NBA, the \texorpdfstring{$\Leftarrow$}{<=} direction}
%----------------------------------------------------------------------
\textbf{Chosen topic:} the proof that every S1S-definable language is
$\w$-regular ($\Leftarrow$).\quad\textbf{Grade:} $30L$.

\begin{enumerate}
  \item Prove briefly the \emph{other} direction ($\Rightarrow$): from a
    Büchi automaton, build an equivalent S1S sentence.
  \item Do you recall a language that is definable by a (nondeterministic)
    BA but \emph{not} by a DBA? [The lecture example: \quoted{finitely
    many $a$}.]
  \item Can you translate that property into LTL? [$FG\,\neg a$.]
  \item Express it also in S1S. [$\exists x\,\forall y\,(x<y \rightarrow
    y\notin Q_a)$.]
  \item We want to build an OBDD for $(x \leftrightarrow y)\wedge R$,
    preferably in canonical form. First fix an order of the variables,
    then draw the diagram.
  \item What is the complexity of building such an OBDD? [Exponential in
    the worst case.]
\end{enumerate}

%----------------------------------------------------------------------
\section{Exam 2 --- Model checking: definition}
%----------------------------------------------------------------------
\textbf{Chosen topic:} the definition of model checking.\quad
\textbf{Grade:} $19$.

\begin{enumerate}
  \item Model checking vs.\ testing.
  \item {}[Given a small Kripke structure on states $q_0,\dots,q_3$:] what
    does $GF(a)$ mean / compute on it?
  \item Give a trace that \emph{is} a model of $GF(a)$ and one that is
    \emph{not}.
  \item Translate the property $GF(a)$ into a Büchi automaton. [This is
    essentially \quoted{infinitely many $a$}, the example seen in the
    lecture.]
  \item Can you construct the complement of [the] automaton?
  \item Explain the acceptance condition (semantics) of a Büchi
    automaton.
\end{enumerate}

%----------------------------------------------------------------------
\section{Exam 3 --- Overview of all the automaton types}
%----------------------------------------------------------------------
\textbf{Chosen topic:} an overview of all the kinds of automaton seen in
the course.\quad\textbf{Grade:} not recorded (the LTL translation in Q3
was marked wrong).

\begin{enumerate}
  \item Prove $\mathrm{DBA}\subseteq\mathrm{DMA}$.
  \item (LTL) Give an example of a model of $GF(a)\leftrightarrow
    GF(b)$.
  \item Translate from natural language into LTL: \quoted{today is sunny,
    but tomorrow it will rain}.
\end{enumerate}

\begin{note}
  For Q3 the student wrote the S1S-style answer $S(x)\wedge R(x{+}1)$,
  which was marked wrong ($\otimes$): LTL has \emph{no} explicit position
  variable~$x$. The intended answer is $\mathit{sunny}\wedge
  X\,\mathit{rain}$ (the examiner's hint was essentially \quoted{what is
  $x$?}).
\end{note}

%----------------------------------------------------------------------
\section{Exam 4 --- Learning automata}
%----------------------------------------------------------------------
\textbf{Chosen topic:} learning automata.\quad\textbf{Grade:} $18$.

\begin{enumerate}
  \item Give a Büchi automaton for \quoted{finitely many $a$}.
  \item $GF(p)\leftrightarrow GF(q)$: what does $GF(p)$ mean? Draw it.
  \item Closure properties of BA ($\cup,\cap,\neg$, and projection).
  \item Prove the case of union.
\end{enumerate}

%----------------------------------------------------------------------
\section{Exam 5 --- The congruence \texorpdfstring{$\approx_{\A}$}{~A}}
%----------------------------------------------------------------------
\textbf{Chosen topic:} \quoted{$\approx_{\A}$ saturates the language, has
finite index, partitions [it]}.\quad\textbf{Grade:} not recorded.

\begin{enumerate}
  \item Prove closure of BA under union.
  \item Give a BA accepting the alternation of $p$ and $q$, i.e.\
    $(pq)^{\w}$. [The scribble writes $(pq)^*$; for an
    $\w$-automaton the intended language is the infinite word $(pq)^{\w}$.]
\end{enumerate}

%----------------------------------------------------------------------
\section{Exam 6 --- LTL}
%----------------------------------------------------------------------
\textbf{Chosen topic:} LTL.\quad\textbf{Grade:} not recorded (Q5 was
marked \emph{failed}).

\begin{enumerate}
  \item Express $G$ and $F$ in basic LTL [i.e.\ using only $X$ and $U$].
  \item BA for \quoted{infinitely many $a$}.
  \item Acceptance condition of a BA. [Side note in the source:
    $GF(p)$ means \quoted{infinitely many $p$}.]
  \item Prove that BA are closed under union --- remember you cannot use
    $\e$-moves.
  \item BA for \quoted{finitely many $a$}. \textbf{--- failed.}
\end{enumerate}

%----------------------------------------------------------------------
\section{Exam 7 --- [topic not recorded]}
%----------------------------------------------------------------------
\textbf{Chosen topic:} not written down; the first question is
\quoted{finitely many $a$}.\quad\textbf{Grade:} $23$.

\begin{enumerate}
  \item \quoted{finitely many $a$}.
  \item How many computations does this automaton have?
  \item S1S and Büchi's theorem: prove the $\Rightarrow$ direction.
  \item Some translations into LTL.
  \item The definition of model checking.
\end{enumerate}

%----------------------------------------------------------------------
\section{Exam 8 --- [topic not recorded: DBA / S1S / OBDD]}
%----------------------------------------------------------------------
\textbf{Chosen topic:} not written down.\quad\textbf{Grade:} not
recorded. (The page is partly the student's worked reasoning.)

\begin{enumerate}
  \item Can we use the union of DBAs to obtain an NBA? [Yes: DBA are
    closed under union --- but one must be careful with the product
    construction.]
  \item Büchi's theorem ($\Leftarrow$), base case $\succxy$. Since the
    successor is injective, $\forall x\in X\,\forall y\in Y.\;x{+}1=y$
    already forces $X$ and $Y$ to be singletons, so one may drop the
    explicit \quoted{$X,Y$ singletons} requirement.
  \item Quantifier (quantification) depth.
  \item OBDD: give the definition and construct one for $(x_1
    \leftrightarrow x_2)\vee x_3$. [Then merge equal outputs and remove
    redundant nodes.]
\end{enumerate}

%----------------------------------------------------------------------
\section{Exam 9 --- DBA}
%----------------------------------------------------------------------
\textbf{Chosen topic:} deterministic Büchi automata.\quad\textbf{Grade:}
not recorded.

\begin{enumerate}
  \item State and prove the theorem: a language is recognised by a DBA
    iff it is the \quoted{limit} --- the vectorial closure
    $\overrightarrow{W}$ --- of a regular set $W$.
  \item \quoted{finitely many $a$}.
  \item Define compassion and justice in LTL.
\end{enumerate}
Further follow-ups noted on the side:
\begin{itemize}
  \item why does refinement \emph{partition} [the language]?
  \item with respect to $\approx_{\A}$;
  \item closure under intersection.
\end{itemize}

%----------------------------------------------------------------------
\section{Exam 10 --- LTL+P succinctness}
%----------------------------------------------------------------------
\textbf{Chosen topic:} LTL with past is (at most exponentially) more
succinct than LTL.\quad\textbf{Grade:} $29$.\quad\textbf{Language:}
Italian.

\medskip
\textit{Warm-up before the chosen topic:} definitions of DBA and NBA and
the DBA$\to$NBA conversion; the definition of LTL and the
NBA$\to$LTL-formula conversion; and an example translating between LTL
and S1S.

\begin{enumerate}
  \item {}[Chosen topic.] There is no LTL+P formula whose equivalent LTL
    formula is only \emph{sub}-exponentially larger (in $n$); i.e.\ LTL+P
    can be exponentially more succinct than LTL.
  \item A Büchi automaton recognising the language over $\mathit{AP}_n =
    \set{p_0,\dots,p_n}$: can one be built in time polynomial in $n$? How
    much memory is needed to store the state at instant $0$? ($n$ bits
    $\Rightarrow$ $2^n$ states\ldots).
  \item In S1S, is the ordering $<$ necessary? No --- it is definable
    from $+1$. Prove it.
  \item Is S1S without $\forall$ as expressive as full S1S? Why may I
    replace each $\forall$ by $\neg\exists$? [And how does this put the
    formula into normal form?]
  \item Formal definition of model checking: what are the inputs and the
    output? [Input: a program, a property, and an optional initial
    state. Output: an exhaustive verdict (this is what distinguishes it
    from testing).]
\end{enumerate}
\begin{note}
  The source notes \quoted{difficulty on the second part of Q4}.
\end{note}

%----------------------------------------------------------------------
\section{Exam 11 --- Closure of BA under intersection}
%----------------------------------------------------------------------
\textbf{Chosen topic:} the proof that Büchi automata are closed under
intersection.\quad\textbf{Grade:} not recorded.\quad\textbf{Language:}
Italian.

\begin{enumerate}
  \item {}[Chosen topic.] Prove closure of BA under intersection. [The
    product construction uses a flag in $\set{1,2,3}$ that tracks which
    automaton's final set was visited; the accepting set is
    $F = Q_1\times Q_2\times\set{3}$, \emph{not}
    $F_1\times F_2\times(\dots)$.]
  \item Why can we \emph{not} proceed via the complementation of a Muller
    automaton?
    \begin{itemize}
      \item define a Muller automaton;
      \item explain how its complementation works.
    \end{itemize}
    Then: which model-checking problems are solvable with automata, in
    particular with Büchi automata? Why is satisfiability done with
    Büchi and not with another type of automaton? [The student did not
    know this second part and left it.]
\end{enumerate}

%----------------------------------------------------------------------
\section{Exam 12 --- S1S \texorpdfstring{$\equiv$}{=} \texorpdfstring{$\w$}{omega}-regular (Büchi's theorem)}
%----------------------------------------------------------------------
\textbf{Chosen topic:} S1S is equivalent to $\w$-regularity.\quad
\textbf{Grade:} $30L$ (\quoted{knew everything}).\quad\textbf{Language:}
Italian.

\begin{enumerate}
  \item {}[Chosen topic.] S1S $\equiv$ $\w$-regular. The direction
    $\mathrm{BA}\subseteq\mathrm{S1S}$ was sketched quickly (mainly the
    acceptance / \quoted{final} relation). The prenex form $\exists Y_0
    \cdots \exists Y_m\, \p(Y_0,\dots,Y_m)$ was introduced --- the same
    normal form sought as the answer to Q4 \quoted{two students ago}
    (see Exam~10, Q4). Then $\text{S1S}_0$: the atoms $X\subseteq Y$ and
    $\succxy$; the connectives $\wedge,\vee,\neg$; and the quantifiers
    $\forall,\exists$, where $\exists$ corresponds to projection.
  \item Why don't we use only Muller automata [in the translation]?
    Because of the existential quantifier (projection): it introduces
    non-determinism.
  \item We saw that DBA are less expressive; can we characterise them?
    [Yes: $\overrightarrow{W}\subseteq\mathrm{DBA}$ and
    $\mathrm{DBA}\subseteq\overrightarrow{W}$ --- the DBA-recognisable
    languages are exactly the vectorial closures of regular sets.]
  \item OBDD: take $\p = (x_1 \leftrightarrow x_2)\wedge x_1$. Explain
    OBDDs briefly and give the canonical OBDD for $\p$ [which simplifies
    to $x_1\wedge x_2$].
\end{enumerate}
\begin{note}
  Remark made during the exam: an OBDD may be drawn roughly and then
  cleaned up afterwards --- one need not get it right on the first try.
\end{note}

%----------------------------------------------------------------------
\section{Exam 13 --- Muller automata \texorpdfstring{$\equiv$}{=} Boolean combinations}
%----------------------------------------------------------------------
\textbf{Chosen topic:} Muller automata recognise exactly the Boolean
(vectorial) combinations of regular sets.\quad\textbf{Grade:} $28$ (fine
through Q3, insufficient on Q4).\quad\textbf{Language:} Italian.\quad
(The source notes this was a \quoted{new topic --- extra points}.)

\begin{enumerate}
  \item {}[Chosen topic.] A language is DMA-recognisable iff it is a
    Boolean combination of vectorial closures $\overrightarrow{W}$ of
    regular sets.
  \item LTL: express the alternation of $p$ and $\neg p$:
    \[
      G\big((p\rightarrow X\neg p)\wedge(\neg p\rightarrow X p)\big).
    \]
  \item Describe the synthesis problem. [Church's problem: realise a
    black-box function from an input stream to an output stream by a
    finite-state (Mealy) machine; solved via Muller automata / games.]
  \item What is the automata-theoretic approach to model checking? What
    operation realises it? [Build the automaton for $\neg\p$, take its
    product / intersection with the system, then test emptiness.]
  \item Explain justice and compassion. Why is justice not a problem for
    LTL (it can be put into a tableau), whereas compassion is?
\end{enumerate}

%----------------------------------------------------------------------
\section{Exam 14 --- Closure of \texorpdfstring{$\w$}{omega}-regular languages under union}
%----------------------------------------------------------------------
\textbf{Chosen topic:} closure of $\w$-regular languages under
union.\quad\textbf{Grade:} not recorded (left blank).\quad
\textbf{Language:} Italian.

\begin{enumerate}
  \item {}[Chosen topic.] Prove that $\w$-regular languages are closed
    under union. --- The student did not know the definition of the
    transition relation $\D\colon Q\times A\to 2^Q$ and left the answer
    blank (\textit{lasciato}).
\end{enumerate}

%----------------------------------------------------------------------
\section{Exam 15 --- Learning automata}
%----------------------------------------------------------------------
\textbf{Chosen topic:} learning automata.\quad\textbf{Grade:} $30L$
(\quoted{knew everything}).\quad\textbf{Language:} Italian.

\begin{enumerate}
  \item {}[Chosen topic.] Is $u\,t =_{L_0} v\,t$ a right invariance (right
    congruence)? [Discuss why / why not.]
  \item Give an example of a language that is \emph{not} $\w$-regular.
    Define an \quoted{ultimately periodic word}. For $\w$-regular $L_1,
    L_2$, relate $\mathrm{UP}(L_1)$ and $\mathrm{UP}(L_2)$ to $L_1,L_2$
    [the characterisation theorem: equal sets of ultimately periodic
    words imply equal languages].
  \item Justice on the (maximal and non-maximal) connected components: for
    all $\t$, either some state has $\t$ disabled, or $\t$ is taken. Why
    is compassion different? For all $\t$, either \emph{no} state has
    $\t$ enabled, or some state has $\t$ taken.
  \item Tamarin: given a pair $(x,y)$, what can the attacker derive about
    $x$ and about $y$? [Under Dolev--Yao the attacker can project a pair
    into its two components; conversely it needs both components to
    construct the pair.]
\end{enumerate}

%======================================================================
\chapter{Questions organised by topic}\label{cha:bytopic}
%======================================================================

The same questions, regrouped by subject, each with a short model answer
and a pointer to the relevant part of the course notes. References such
as \quoted{(Exam~6, Q1)} point back to \autoref{cha:bysession}.

%----------------------------------------------------------------------
\section{Büchi automata and \texorpdfstring{$\w$}{omega}-regular languages}
%----------------------------------------------------------------------

\subsection{\quoted{Finitely many \texorpdfstring{$a$}{a}} and \quoted{infinitely many \texorpdfstring{$a$}{a}}}
\emph{Asked in Exams 1, 4, 6, 7, 9 (and as an S1S/LTL translation in
Exam~1).}

These two languages are the workhorses of the course because they
separate deterministic from nondeterministic Büchi automata.

\paragraph{Infinitely many $a$ ($\exists^{\w}n.\,\a(n)=a$).}
DBA-recognisable. Two states $q_0,q_1$; on $a$ go to the final state
$q_1$, on $b$ go back to $q_0$; $F=\set{q_1}$. A run is accepting iff it
visits $q_1$ infinitely often, i.e.\ reads $a$ infinitely often. In LTL:
$GF\,a$. (This is a \emph{reset automaton}: each symbol leads to a single
state regardless of the current one.)

\paragraph{Finitely many $a$ ($\exists^{<\w}n.\,\a(n)=a$).}
$\w$-regular but \textbf{not} DBA-recognisable. A (nondeterministic) BA
\quoted{guesses} the last position of an $a$: stay in $q_0$ reading
anything, then nondeterministically move to a final state $q_1$ from
which only $b$ is allowed forever; $F=\set{q_1}$.
\begin{itemize}
  \item \textbf{LTL:} $FG\,\neg a$ --- eventually, always $\neg a$.
  \item \textbf{S1S:} $\exists x\,\forall y\,(x<y \rightarrow y\notin
    Q_a)$ --- from some position on, no symbol is an $a$.
  \item \textbf{Why no DBA?} If a DBA recognised it as
    $\overrightarrow{W}$ for some regular $W$, one could splice
    arbitrarily many $a$'s back in (because $b^{\w}$ has infinitely many
    prefixes in $W$, then $b^{n_1}a b^{\w}$ does too, etc.) and build a
    word with infinitely many $a$'s that the automaton would still
    accept --- contradiction.
\end{itemize}
Failing to produce the BA for \quoted{finitely many $a$} is exactly what
sank Q5 in Exam~6: the trap is to try to do it deterministically.

\subsection{Acceptance condition; the meaning of \texorpdfstring{$GF$}{GF} and \texorpdfstring{$FG$}{FG}}
\emph{Asked in Exams 2, 4, 6.}

A computation $\s$ of a BA $\A=\NBA$ on an $\w$-word $\a$ starts in $q_0$,
respects $\D$ at every step, and is \textbf{successful} iff
$\infs\cap F\neq\emptyset$ --- i.e.\ it visits some final state
\emph{infinitely often}. $\A$ accepts $\a$ iff \emph{some} computation on
$\a$ is successful (note \quoted{some}, not \quoted{the}: there are
in general infinitely many computations). Correspondingly, as an LTL
slogan:
\[
  GF\,p \;=\; \text{$p$ holds infinitely often},\qquad
  FG\,p \;=\; \text{$p$ eventually holds forever (stability)}.
\]

\subsection{How many computations does an automaton have? (Exam 7, Q2)}
On an infinite input a \emph{nondeterministic} automaton has in general
infinitely (even uncountably) many computations --- at each step every
choice allowed by $\D$ spawns a branch. A \emph{deterministic} one has
exactly one computation per input word. This is precisely why acceptance
is defined existentially (\quoted{there exists a successful
computation}).

\subsection{Closure properties and the \texorpdfstring{$\e$}{epsilon}-move caveat}
\emph{Asked in Exams 4, 5, 6, 8, 11, 14.}

$\w$-regular languages are closed under $\cup$, $\cap$, complementation,
left concatenation with a regular set, and $\w$-iteration.

\paragraph{Union (Exams 4, 5, 6, 14).}
Take the two BA on disjoint state sets. Unlike for NFAs, you \emph{cannot}
just add a fresh initial state with $\e$-moves to the two old initial
states (Büchi automata as defined here have no $\e$-transitions). The
clean fix is a nondeterministic initial choice: duplicate the outgoing
transitions of both old initial states onto a single new initial state
(and keep it final if either old one was). This is the very point the
examiner stressed in Exam~6, Q4.

\paragraph{Intersection (Exam 11).}
A simple product is not enough, because the two automata may visit their
final states at different times. The standard construction takes
$Q_1\times Q_2\times\set{1,2,3}$ with a flag that advances $1\to2$ when
$\A_1$ hits $F_1$, $2\to3$ when $\A_2$ hits $F_2$, then resets $3\to1$;
the accepting set is
\[
  F \;=\; Q_1\times Q_2\times\set{3},
\]
\emph{not} $F_1\times F_2\times(\dots)$. Reaching flag $3$ infinitely
often certifies that \emph{both} automata accept.

\paragraph{Why projection/$\Delta$ matters (Exam 14).}
The student in Exam~14 could not state $\D\colon Q\times A\to 2^Q$ --- the
nondeterministic transition relation. It is the backbone of every BA
construction (union, projection, the S1S translation), which is why not
knowing it blocks the whole proof.

\subsection{Complementation and the congruence \texorpdfstring{$\approx_{\A}$}{~A}}
\emph{Chosen topic of Exam 5.}

Given a BA $\A$, define $u\approx_{\A}v$ iff for all states $s,s'$:
$s\to_u s' \iff s\to_v s'$ and $s\to_u^F s' \iff s\to_v^F s'$ (the second
relation additionally requires a final state along the way). Then
$\approx_{\A}$ is a \textbf{congruence of finite index} that
\textbf{saturates} $\L(\A)$:
\[
  UV^{\w}\cap L\neq\emptyset \;\Longrightarrow\; UV^{\w}\subseteq L .
\]
Saturation of $L$ implies saturation of $\overline{L}$, and every
$\w$-word lies in some $UV^{\w}$ with $VV\subseteq V$. Hence both $L$ and
$\overline{L}$ are finite unions $\bigcup_i U_iV_i^{\w}$ of
$\w$-regular pieces, giving closure under complementation
\emph{effectively}. This is the machinery behind \quoted{$\approx_{\A}$
saturates, has finite index, partitions}.

\subsection{Emptiness, \texorpdfstring{$\w$}{omega}-regular expressions and ultimately periodic words}
\emph{Asked in Exam 15.}

\begin{itemize}
  \item \textbf{Emptiness} of a BA is decidable and easy: it is a
    reachability problem --- reach a final state from $q_0$ and find a
    cycle through it.
  \item An \textbf{$\w$-regular expression} has the form
    $\bigcup_{i=1}^n U_i\cdot V_i^{\w}$ with $U_i,V_i$ regular; BA and
    $\w$-regular expressions are equi-expressive.
  \item An \textbf{ultimately periodic} word is $u\,v^{\w}$ with
    $u,v\in A^*$. Every non-empty $\w$-regular language contains one (so
    a language with \emph{no} ultimately periodic word --- e.g.\ the
    common-suffix language of a non-ultimately-periodic $\b$ --- is
    \textbf{not} $\w$-regular; this answers the \quoted{non-$\w$-regular
    example} of Exam~15).
  \item \textbf{Characterisation (Exam 15, Q2):} an $\w$-regular language
    is uniquely determined by its set $\mathrm{UP}(L)$ of ultimately
    periodic words. Hence $\mathrm{UP}(L_1)=\mathrm{UP}(L_2)\Rightarrow
    L_1=L_2$ (contrapositive: $L_1\neq L_2$ forces a difference already
    among ultimately periodic words, because $L_1\setminus L_2$ is
    $\w$-regular and non-empty, hence contains one).
\end{itemize}

%----------------------------------------------------------------------
\section{DBA, Muller, Rabin and the expressiveness hierarchy}
%----------------------------------------------------------------------
\emph{Asked in Exams 3, 8, 9, 11, 12, 13.}

\subsection{\texorpdfstring{$\mathrm{DBA}\subseteq\mathrm{DMA}$}{DBA <= DMA} (Exam 3, Q1)}
Given a DBA $\DBA$, the equivalent DMA keeps the same transition
structure and sets $\F=\set{P\in 2^Q : P\cap F\neq\emptyset}$. A Muller
computation is successful iff the set of infinitely-often-visited states
lies in $\F$, i.e.\ contains a final state --- exactly the Büchi
condition.

\subsection{Characterising DBA via the vectorial closure (Exams 9, 12)}
$L$ is DBA-recognisable iff $L=\overrightarrow{W}$ for a regular $W$,
where $\overrightarrow{W}=\set{\a : \a \text{ has infinitely many
prefixes in } W}$. Both inclusions ($\overrightarrow{W}\subseteq
\mathrm{DBA}$ and $\mathrm{DBA}\subseteq\overrightarrow{W}$) are the
content of Exam~9 Q1 and Exam~12 Q3. Since DBA are not closed under
complementation, $\mathrm{DBA}\subsetneq\mathrm{BA}$.

\subsection{Muller \texorpdfstring{$\equiv$}{=} Boolean combinations of vectorial closures (Exams 11, 13)}
A language is DMA-recognisable iff it is a Boolean combination (union,
intersection, complement) of languages $\overrightarrow{W}$ with $W$
regular. Together with McNaughton's theorem
($\mathrm{NBA}\equiv\mathrm{DMA}$) this yields the final hierarchy
\[
  \mathrm{DBA}\;\subsetneq\;\mathrm{DMA}\equiv\mathrm{DRA}\equiv
  \mathrm{NBA}\equiv\mathrm{NMA}\equiv\mathrm{NRA},
\]
i.e.\ everything except DBA is equally expressive. The
\quoted{complementation of a Muller automaton} question (Exam~11, Q2)
turns on this: complementation is easy for the \emph{deterministic}
Muller/Rabin condition but not for nondeterministic Büchi directly.

\subsection{Union of DBAs to get an NBA (Exam 8, Q1)}
Yes: DBA are closed under union (and intersection) via the product
construction --- but the construction goes through products of states,
and the result is read as a Büchi/Muller condition; one cannot naively
merge initial states.

%----------------------------------------------------------------------
\section{S1S and Büchi's theorem}
%----------------------------------------------------------------------
\emph{Asked in Exams 1, 7, 8, 10, 12.}

\subsection{The two directions}
\textbf{$\Rightarrow$ (BA $\to$ S1S; Exams 1 Q1, 7 Q3, 12 Q1).}
From $\A=\NBA$ with $Q=\set{0,\dots,m}$, write $\exists Y_0\cdots\exists
Y_m\,\p(Y_0,\dots,Y_m)$ where $Y_i$ collects the positions in state $i$.
The matrix $\p$ has three conjuncts: start in $q_0$ with a unique state
per position; respect $\D$ at every step; and visit some final state
infinitely often (the \quoted{final relation}
$\bigvee_{i\in F}\forall x\,\exists y\,(x<y\wedge y\in Y_i)$).

\textbf{$\Leftarrow$ (S1S $\to$ BA; chosen in Exam 1).}
Done by induction on formula structure after reducing to $\text{S1S}_0$.

\subsection{\texorpdfstring{$\text{S1S}_0$}{S1S0} and the successor base case (Exams 8, 12)}
$\text{S1S}_0$ keeps only second-order variables, with atoms $X\subseteq
Y$ and $\succxy$ (meaning $X=\set{x},Y=\set{y},x{+}1=y$). The translation
is then inductive: $\wedge,\vee$ by intersection/union, $\neg$ by
complementation, and $\exists$ by \textbf{projection} (Exam~12, Q1). The
base automata for $X\subseteq Y$ and $\succxy$ are trivial. The remark in
Exam~8 Q2 is that, since the successor is injective, $\forall x\in
X\,\forall y\in Y.\,x{+}1=y$ already pins $X,Y$ to singletons.

\subsection{Why nondeterminism / Muller (Exam 12, Q2)}
Projection ($\exists X$) is implemented by \emph{erasing} one track of
the transition labels, which is sound precisely because BA are
nondeterministic. A purely deterministic device (e.g.\ deterministic
Muller alone) cannot perform projection directly; hence the construction
must pass through nondeterministic Büchi automata.

\subsection{\texorpdfstring{$<$}{<} from \texorpdfstring{$+1$}{+1}, \texorpdfstring{$\forall$}{forall} from \texorpdfstring{$\neg\exists$}{not-exists}, quantifier depth (Exams 8, 10)}
\begin{itemize}
  \item \textbf{$<$ is definable from $+1$} (Exam~10, Q3):
    $x<y \equiv \forall X\,[\,x{+}1\in X \wedge \forall z(z\in
    X\rightarrow z{+}1\in X)\rightarrow y\in X\,]$. The converse needs
    second order, so $<$ is strictly more expressive than $+1$ at first
    order.
  \item \textbf{$\forall$ as $\neg\exists$} (Exam~10, Q4): $\forall
    x\,\p\equiv\neg\exists x\,\neg\p$. Rewriting a prenex formula this
    way turns it into alternating blocks $\exists\cdots\neg\exists\cdots$;
    the number of negations equals the number of quantifier alternations.
  \item \textbf{Quantifier depth} (Exam~8, Q3) is the maximum nesting of
    quantifiers. It controls everything: membership in the S1S theory is
    \emph{non-elementary} (a tower of exponentials whose height is the
    alternation depth), even though BA emptiness is polynomial --- the
    blow-up lives in the formula-to-automaton translation. (WS1S, the
    weak variant, is the version Büchi actually used; it is equivalent to
    S1S.)
\end{itemize}

%----------------------------------------------------------------------
\section{Linear Temporal Logic (LTL)}
%----------------------------------------------------------------------
\emph{Asked in Exams 2, 3, 4, 6, 13 (and translations in 1, 7).}

\subsection{Deriving \texorpdfstring{$G$}{G} and \texorpdfstring{$F$}{F} from the basic operators (Exam 6, Q1)}
The primitive modalities are $X$ (next) and $U$ (until). Then
\[
  F\,\p \;\equiv\; \top\,U\,\p, \qquad G\,\p \;\equiv\; \neg F\,\neg\p .
\]

\subsection{Natural language to LTL, and \quoted{what is \texorpdfstring{$x$}{x}?} (Exam 3, Q3)}
\quoted{Today is sunny, but tomorrow it will rain} becomes
\[
  \mathit{sunny}\wedge X\,\mathit{rain}.
\]
The wrong answer given in the exam, $S(x)\wedge R(x{+}1)$, is
\emph{S1S} thinking: LTL is evaluated at an implicit \quoted{current
instant} and has no position variable $x$ to quantify or index. Time
advances only through the modalities ($X$, $U$, $F$, $G$).

\subsection{Alternation of \texorpdfstring{$p$}{p} and \texorpdfstring{$\neg p$}{not p} (Exam 13, Q2)}
\[
  G\big((p\rightarrow X\neg p)\wedge(\neg p\rightarrow X p)\big),
\]
i.e.\ the truth value of $p$ flips at every step.

\subsection{Models, non-models and \texorpdfstring{$GF(a)\leftrightarrow GF(b)$}{GF(a)<->GF(b)}}
\emph{Exams 2, 4 (meaning and models), 3 (the biconditional).}
$GF\,a$ holds on a trace iff $a$ is true at infinitely many positions; so
$(\set{a,b})^{\w}$ and $(\set{b}\set{b}\set{a}\set{b})^{\w}$ satisfy it,
whereas $\set{a}\cdot(\set{b})^{\w}$ does not. A model of
$GF(a)\leftrightarrow GF(b)$ (Exam~3, Q2) is any trace where $a$ and $b$
are \emph{both} infinitely frequent (e.g.\ $(\set{a,b})^{\w}$) or
\emph{both} only finitely frequent (e.g.\ $\emptyset^{\w}$); a trace with
infinitely many $a$ but finitely many $b$ falsifies it.

\subsection{Kamp's theorem}
LTL captures exactly the \emph{first-order} fragment of S1S over linear
orders. This frames every \quoted{translate to LTL / to S1S} exercise:
LTL = FO over $(\NN,<)$, while full S1S adds (monadic) second-order
power.

%----------------------------------------------------------------------
\section{LTL with past, and succinctness (Exam 10)}
%----------------------------------------------------------------------
LTL+P adds the past modalities $Y$ (yesterday) and $S$ (since), with
derived $O$ (once), $H$ (historically) and weak yesterday $\widetilde{Y}$.

\begin{itemize}
  \item \textbf{Same expressiveness:} LTL+P is exactly as expressive as
    LTL (past adds no new languages).
  \item \textbf{Strictly more succinct (the chosen topic):} LTL+P can be
    \emph{exponentially} more succinct. The witness family $L_n$ over
    $\set{p_0,\dots,p_n}$ has a linear-size LTL+P formula but every LTL
    formula for it is at least exponential --- because every NBA for
    $L_n$ is doubly exponential ($2^{2^{\mathcal{O}(n)}}$) while each LTL/LTL+P
    formula yields an at-most-exponential NBA. Exam~10 Q2's
    \quoted{store the state at instant $0$ in $n$ bits $\Rightarrow 2^n$
    states} is the intuition for that lower bound.
\end{itemize}

%----------------------------------------------------------------------
\section{Fairness: justice vs.\ compassion}
%----------------------------------------------------------------------
\emph{Asked in Exams 9, 13, 15.}

\textbf{Justice} (weak fairness): a transition that is \emph{continuously}
enabled cannot be ignored forever. \textbf{Compassion} (strong fairness):
a transition that is enabled \emph{infinitely often} cannot be taken only
finitely often. Compassion implies justice, not conversely.

\paragraph{In LTL (Exams 9 Q3, 13 Q5).}
\[
  \text{justice:}\quad GF(\mathit{enabled})\;\text{eventually settled}\;
  \rightarrow GF(\mathit{taken}),\qquad
  \text{compassion:}\quad GF(\mathit{enabled})\rightarrow GF(\mathit{taken}).
\]
More precisely justice is $\neg F\big(G(\mathit{enabled})\wedge
G(\neg\mathit{taken})\big)$ and compassion is
$\neg\big(GF(\mathit{enabled})\wedge FG(\neg\mathit{taken})\big)$. Note
$GF(\mathit{enabled}\rightarrow\mathit{taken})$ is \emph{not} the same ---
it would force \quoted{taken at the same instant as enabled}.

\paragraph{On components / supergraphs (Exams 13 Q5, 15 Q3).}
For a strongly connected subgraph $S$:
\begin{itemize}
  \item $S$ is \emph{just} if every $\t\in\J$ is either taken in $S$ or
    disabled at \emph{some} node of $S$;
  \item $S$ is \emph{compassionate} if every $\t\in\C$ is either taken in
    $S$ or disabled at \emph{all} nodes of $S$.
\end{itemize}
Justice is preserved by enlarging the graph (it holds for supergraphs),
so it can be handled directly in the tableau --- this is why
\quoted{justice is not a problem for LTL}. Compassion is \emph{not}
preserved by supergraphs (a larger SCS may enable a transition that the
smaller one never did), so MSCS-checking alone is insufficient and a
dedicated decomposition is needed --- this is why \quoted{compassion is
different}.

%----------------------------------------------------------------------
\section{Model checking}
%----------------------------------------------------------------------
\emph{Asked in Exams 2, 7, 10, 13.}

\subsection{Definition, inputs and output (Exams 2, 7, 10)}
Given a Kripke structure $M$ (the system), a state $s$, and a temporal
formula $\p$, model checking decides whether $M,s\models A\p$ (all paths
from $s$ satisfy $\p$). Inputs: a model/program, a property, and
(optionally) an initial state; output: a yes/no verdict and, on failure,
a \textbf{counterexample} trace. LTL model checking is PSPACE-complete.

\subsection{Model checking vs.\ testing (Exam 2, Q1)}
Model checking is automatic, \textbf{exhaustive}, and unbiased (it does
not privilege \quoted{likely} scenarios), and it returns a counterexample
when a property fails. Testing only samples behaviours. The price is the
\emph{state-space explosion} and dependence on the quality of the model.

\subsection{The automata-theoretic approach (Exam 13, Q4)}
Build a Büchi automaton $\A_{\neg\p}$ for the negation of the property,
take its product (intersection) with the system automaton, and check the
result for \textbf{emptiness}: $M\models\p$ iff $\L(M)\cap\L(\A_{\neg\p})
=\emptyset$. The key \quoted{operation} is therefore
intersection-then-emptiness. (Equivalently, in the tableau/behaviour-graph
formulation of the notes: build $\B_{(P,\p)}$ and search for an adequate
--- fair and fulfilling --- subgraph.)

%----------------------------------------------------------------------
\section{Symbolic methods and OBDDs}
%----------------------------------------------------------------------
\emph{Asked in Exams 1, 8, 12.}

\subsection{What an OBDD is}
An Ordered Binary Decision Diagram is a reduced, ordered decision graph
for a Boolean function. Once a variable order is fixed, the reduced OBDD
is a \textbf{canonical form}: it is unique and minimal, so equivalence of
formulas reduces to graph isomorphism. Reduction applies three rules
until stable: (1) merge equal leaves; (2) merge isomorphic internal nodes
(same variable, same $\low$, same $\high$); (3) delete redundant nodes
($\low(v)=\high(v)$). The remark from Exam~12 --- draw it roughly, then
clean up --- is exactly this reduction process.

\subsection{The worked formulas}
\begin{itemize}
  \item \textbf{$(x\leftrightarrow y)\wedge R$} (Exam 1, Q5), order
    $x,y,R$: test $x$; on each branch test $y$ and keep going only if
    $y=x$ (else $\to 0$); finally test $R$ ($R{=}1\to 1$, $R{=}0\to 0$).
    The two \quoted{$y$ matches $x$} branches share the same $R$-node,
    which merges in the reduced diagram.
  \item \textbf{$(x_1\leftrightarrow x_2)\wedge x_1$} (Exam 12, Q4)
    simplifies to $x_1\wedge x_2$: test $x_1$ ($0\to0$), then $x_2$
    ($0\to0$, $1\to1$). A clean two-node OBDD.
  \item \textbf{$(x_1\leftrightarrow x_2)\vee x_3$} (Exam 8, Q4): test
    $x_1,x_2$; if $x_1=x_2$ the function is already $1$; otherwise it
    equals $x_3$. So the \quoted{$x_1=x_2$} paths go to $1$ and the
    \quoted{$x_1\neq x_2$} paths go to an $x_3$-node.
\end{itemize}

\subsection{Complexity (Exam 1, Q6)}
In the worst case an OBDD is \textbf{exponential} in the number of
variables, and finding the optimal variable order is NP-complete; good
heuristics (keep variables that interact close together) usually keep it
manageable. Symbolically, $\exists x\,f \equiv f|_{x\leftarrow 0}\vee
f|_{x\leftarrow 1}$ and $\forall x\,f\equiv f|_{x\leftarrow 0}\wedge
f|_{x\leftarrow 1}$ (the \textsc{Restrict}/\textsc{Apply} operations),
which is what makes symbolic model checking over OBDDs work.

%----------------------------------------------------------------------
\section{The synthesis problem (Exam 13, Q3)}
%----------------------------------------------------------------------
Church's synthesis problem: given a specification $\p(X,Y)$ relating an
\emph{input} stream $X$ (chosen by the environment) to an \emph{output}
stream $Y$ (chosen by us), build a finite-state (Mealy) machine that
reads $X$ bit-by-bit and produces $Y$ on the fly so that $(\a,\b)\models
\p$ --- or report that none exists. It is the realizability (not just
satisfiability) question, naturally read as an infinite two-player game:
the environment plays inputs, the system replies with outputs, and the
system wins iff the specification holds. The Büchi--Landweber theorem
solves it for S1S/Muller specifications: every Muller game on a finite
arena is determined with finite-state strategies (memory bounded by
$n!\cdot n$ via the Latest Appearance Record). The exam's phrasing ---
\quoted{a black-box function input$\to$output, realised with Muller
automata} --- is this setup.

%----------------------------------------------------------------------
\section{Protocol verification: Tamarin (Exam 15, Q4)}
%----------------------------------------------------------------------
Tamarin verifies security protocols in the symbolic
\textbf{Dolev--Yao} model, where the attacker controls the network: it
can read, intercept, block, replay and \emph{forge} messages from its
current knowledge. Knowledge is closed under the equational theory and
the term constructors; in particular \textbf{pairing} is transparent to
the attacker:
\[
  \text{from } \langle x,y\rangle \text{ the attacker derives both } x
  \text{ and } y \text{ (projection)},
\]
and conversely it can build $\langle x,y\rangle$ only if it already knows
\emph{both} $x$ and $y$. (Contrast with encryption: from $\mathtt{enc}(m,
k)$ it gets $m$ only if it also knows $k$.) That is the expected answer
to \quoted{given $(x,y)$, what can the attacker derive about $x$ and
$y$?}

%----------------------------------------------------------------------
\section{Learning automata}
%----------------------------------------------------------------------
\emph{Chosen topic of Exams 4 and 15.}

The learner discovers an unknown regular language $L_0$ by
\textbf{membership} queries (is $w\in L_0$?) and \textbf{equivalence}
queries (is my hypothesis correct? --- answered with a counterexample).
Membership alone is insufficient (infinitely many words); equivalence
alone is only exponential; together (Angluin's $L^*$) the minimal DFA is
learned in polynomially many queries. The hypothesis is built from an
observation table over sets $S$ (representatives) and $T$ (test
suffixes), kept $T$-minimal and $T$-complete.

\paragraph{Right invariance of the relativised relation (Exam 15, Q1).}
The full Myhill--Nerode relation $u\sim_{L_0}v$ (i.e.\ $\forall t.\ ut\in
L_0\iff vt\in L_0$) \emph{is} a right congruence: if $u\sim_{L_0}v$ then
$ua\sim_{L_0}va$, by absorbing the leading $a$ into the test word
($t=at'$). But the \emph{relativised} relation
$\approx_{L_0,T}$ (testing only over a finite $T$) is \textbf{not} always
right invariant: appending a symbol can send two $T$-equivalent words
into different classes, because $T$ may not contain the suffixes that
distinguish them. This asymmetry --- the exam's question --- is exactly
what drives the $L^*$ loop: a counterexample adds the missing suffixes to
$T$, restoring progress toward the (right-invariant, finite-index) limit
$\sim_{L_0}$.

\end{document}
